%=====================================================================
\section{Cyclic operation}
\label{sec:cyclic}

Everything in \S\ref{sec:hydraulics} is a \emph{first-cycle} result. The store
begins at $E'=0$ --- fully solid at exactly $\Tm$, with no subcooling --- and
the cycle does not return there. A plant repeating a daily schedule never sees
that state again after the first day, so the well counts reported above are
cold-start numbers.

The intended duty is \textbf{\SI{10}{\hour} charge, \SI{2}{\hour} rest,
\SI{10}{\hour} discharge, \SI{2}{\hour} rest}.

%---------------------------------------------------------------------
\subsection{The idle periods are exact no-ops}
\label{sec:cyclic:rest}

With no flow, $\mathrm{NTU}\to\infty$ but the conductance of
Eq.~\eqref{eq:K-derived} carries $\md$ as a prefactor:
\begin{equation}
  K = \frac{\md c_{p}\left(1-e^{-\mathrm{NTU}}\right)}{\Delta z}
  \;\xrightarrow[\;\md\to0\;]{}\; 0 ,
\end{equation}
so $\mathrm{d}E'/\mathrm{d}t = K(T_{j}-\Tpcm) = 0$ identically. Hypotheses H2
(no axial conduction), H6 (adiabatic wall) and H9 (no azimuthal exchange) leave
no other path by which a segment can exchange heat with anything. The
\SI{2}{\hour} rests therefore change nothing: the 10/2/10/2 schedule is
arithmetically identical to 10/10 in this model.

That is worth reading as a diagnostic rather than a convenience. Physically, a
\SI{2}{\hour} rest lets superheated liquid near the tube conduct outward into
colder PCM and lets the front relax; the lumped state of H5 cannot represent
that either. The rests doing nothing is a statement about the strength of four
assumptions, not about the store.

%---------------------------------------------------------------------
\subsection{Why sizing cannot simply be moved to cyclic steady state}
\label{sec:cyclic:degenerate}

At cyclic steady state (CSS) the state returns to itself, so the enthalpy change
over a complete cycle is zero. The model has no loss path, and therefore
\begin{equation}
  \boxed{\;\eta_{\mathrm{storage}}(\mathrm{CSS}) \;\equiv\; 1\;}
  \label{eq:css-unity}
\end{equation}
exactly --- for every well count and every flow. Computed values agree to eight
decimal places across $\Nwells$ from $11$ to $20$ and flow ratios from $0.92$ to
$1.03$. Equation~\eqref{eq:css-unity} is not a result about the store; it is
energy conservation applied to a closed cycle in an adiabatic model, and it is
worth stating only because of the two consequences that follow.

\subsubsection*{Consequence 1: the loss surplus is unattainable}

The budget of Eq.~\eqref{eq:budget} requires
$\Delta E_{\mathrm{out,HP}}=(1+\lambda)\Delta E_{\mathrm{in,ORC}}$. With
Eq.~\eqref{eq:css-unity} the ratio of CSS charge to requirement pins at
\begin{equation}
  \frac{\Delta E_{\mathrm{charge}}^{\mathrm{CSS}}}
       {\Delta E_{\mathrm{out,HP}}} \;=\; \frac{1}{1+\lambda}
  \;=\; 0.95238
\end{equation}
for \emph{every} $\Nwells$. The sweep returns $0.95247$ at $\Nwells=12$ and
$0.95231$ at $\Nwells=20$: flat. There is no root, and the CSS sizing problem as
posed in \S\ref{sec:hydraulics:sizing} has no solution at any well count.

It is worth being precise about where the surplus goes on the first cycle, since
nothing in the model loses it. $\lambda$ enters in exactly one place --- the
line $\Delta E_{\mathrm{out,HP}} = (1+\lambda)\Delta E_{\mathrm{in,ORC}}$ ---
and propagates only to the charge target, the charging mass flow, and
$\dot W_{\mathrm{el,in}}$. Nothing checks that the surplus is lost. It is a
requirement to \emph{deposit} more, and on the first cycle the extra is simply
\textbf{parked in the store} as residual melt:

\begin{center}
\begin{tabular}{@{}l r@{}}
\toprule
 & kJ, field total \\
\midrule
charged, actually delivered & \num{1.97546e8} \\
discharged                  & \num{1.88298e8} \\
\midrule
difference                  & \num{9.24816e6} \\
$\lambda\,\Delta E_{\mathrm{in,ORC}}$ & \num{9.41140e6} \\
ratio                       & 0.983 \\
\bottomrule
\end{tabular}
\end{center}

\noindent
The first cycle gets away with it only because it starts from an empty store. At
CSS there is nowhere left to park it.

\medskip
\noindent\textbf{A correction to \S\ref{sec:hydraulics:ideal} as written at
v0.4b.} That section reported $\eta_{\mathrm{storage}}=0.9532$ against the
assumed $1/(1+\lambda)=0.9524$ and offered the agreement as independent
corroboration of the capacity correction. It is not independent. The discharge
loop solves so that the delivered energy equals $\Delta E_{\mathrm{in,ORC}}$,
and the charge is sized so the absorbed energy equals
$(1+\lambda)\Delta E_{\mathrm{in,ORC}}$; therefore
$\eta_{\mathrm{storage}}\equiv 1/(1+\lambda)$ identically whenever both loops
converge. It is a tautology and should not have been offered as evidence.

\subsubsection*{Consequence 2: with $\lambda=0$ the rate criterion goes degenerate}

Setting $\lambda=0$ removes the contradiction --- the target becomes attainable
--- but it does not restore the ability to size. By
Eq.~\eqref{eq:css-unity} the CSS charge and the CSS discharge are then the
\emph{same number}, so the two sizing equations collapse into one. It fixes the
\emph{flow ratio} and says nothing whatever about $\Nwells$:

\begin{center}
\begin{tabular}{@{}r r r r@{}}
\toprule
$\Nwells$ & flow ratio & charge/req & discharge/req \\
\midrule
11.0 & 1.0740 & 1.00000 & 1.00000 \\
14.0 & 1.0280 & 0.99994 & 0.99994 \\
20.0 & 0.9965 & 0.99994 & 0.99994 \\
\bottomrule
\end{tabular}
\end{center}

\noindent
This is the same structural failure that \S\ref{sec:hydraulics:sizing} records
for the discharge side, now appearing on the charge side as well. Cycling is
what exposes it. Note also that adding a loss mechanism --- relaxing H6 ---
would remove the \emph{contradiction} of Consequence~1 but not this degeneracy:
at CSS the charge would then exceed the discharge by exactly the loss, and the
one remaining equation would still pin the flow ratio rather than $\Nwells$.

%---------------------------------------------------------------------
\subsection{The reformulation: specify, simulate, report}
\label{sec:cyclic:reform}

The resolution is to stop asking the energy balance to determine $\Nwells$.
Specify the hardware, march to CSS, and report what comes out:

\begin{center}
\begin{tabular}{@{}l l@{}}
\toprule
quantity & how it is set \\
\midrule
$\Nwells$ & $\Delta E_{\mathrm{out,HP}}/(\rhom V_{\mathrm{well}}\hm)$, latent heat alone \\
$\md^{\mathrm{ch}}$ & pinned by the charging glide \\
$\md^{\mathrm{dc}}$ & pinned by the discharging glide \\
\bottomrule
\end{tabular}
\end{center}

\noindent
Nothing is solved. There is no root-find anywhere in \code{simulate\_css}, so
the question of convergence does not arise; what was an unsatisfiable constraint
becomes a reported output --- the deviation of the delivered energy, and hence
of the net electrical output, from its target. Sizing on latent heat alone
\emph{overestimates} $\Nwells$, because it ignores the sensible capacity the
store also has; that is deliberate and conservative.

This framing also dissolves an older objection. \S\ref{sec:hydraulics:sizing}
argues that the discharge flow cannot be pinned to its glide, because the
delivered energy then saturates below the requirement for any well count and
leaves no root. Under the reformulation that saturation is not a failure: it
\emph{is} the answer.

\subsubsection*{Result at the design point}

\begin{center}
\begin{tabular}{@{}l r l@{}}
\toprule
$\Nwells$ & 11.5441 & latent heat alone, not solved \\
$\md^{\mathrm{ch}}$, $\md^{\mathrm{dc}}$ & 0.9655, 0.9698 & \si{\kilogram\per\second} per leg-pair \\
flow ratio & 1.0044 & a consequence, not a solve \\
cycles to CSS & 28 & \\
$\eta_{\mathrm{storage}}$ & 1.000000 & Eq.~\eqref{eq:css-unity} \\
\midrule
required $\Delta E_{\mathrm{in,ORC}}$ & \num{1.88228e8} & kJ \\
delivered at CSS & \num{1.77417e8} & kJ \\
\textbf{deviation} & \textbf{-5.736} & \si{\percent} \\
\midrule
implied net output & 0.9426 & MWe against a target of 1.0000 \\
\bottomrule
\end{tabular}
\end{center}

\subsubsection*{The shortfall is a rate limit, not an inventory limit}

At CSS the store melts \emph{completely} --- $\varepsilon_{\mathrm{local}}=1.000$
at every segment at end of charge --- but returns only to a mean of $0.154$. It
never fully refreezes within the \SI{10}{\hour} discharge window, and that
residual is what limits delivery. The design curve makes the diagnosis explicit:

\begin{center}
\begin{tabular}{@{}r r r r r@{}}
\toprule
$\Nwells$ & delivered/req & MWe & $\varepsilon$ end charge & $\varepsilon$ end discharge \\
\midrule
11.544 & 0.9426 & 0.9426 & 1.000 & 0.154 \\
12.500 & 0.9540 & 0.9540 & 1.000 & 0.209 \\
14.000 & 0.9662 & 0.9662 & 1.000 & 0.288 \\
16.000 & 0.9777 & 0.9777 & 1.000 & 0.376 \\
20.000 & 0.9931 & 0.9931 & 1.000 & 0.505 \\
26.000 & 1.0082 & 1.0082 & 1.000 & 0.629 \\
\bottomrule
\end{tabular}
\end{center}

\noindent
The residual melt fraction \emph{rises} with $\Nwells$. That is the signature of
a rate-limited discharge: more wells means each is worked less hard, so
proportionally less of each one refreezes. Were the limit an inventory limit,
the residual would fall. The lever is therefore on the discharge side --- flow,
window, or conductance --- and not more PCM. The target is met near
$\Nwells\approx25$, about $2.2\times$ the latent-only count.

\subsubsection*{How much of this to believe}

The \SI{-5.7}{\percent} deviation is a model result and inherits every
assumption of \S\ref{sec:defects}. H9 is worth of order \SI{25}{\percent} of the
useful duty at the wellhead (\S\ref{sec:defects:h9}); H3 is saturated over
\SI{99}{\percent} of the well (\S\ref{sec:defects:h3}); the time grid carries a
one-sided \SI{-0.3}{\percent} bias. A \SI{6}{\percent} deviation sits inside
that band. It should be reported as \emph{``the store is discharge-rate limited
by roughly \SI{6}{\percent} at this sizing''} and not as a calibrated shortfall.

%---------------------------------------------------------------------
\subsection{A defect that cycling exposed}
\label{sec:cyclic:reversal}

The discharge marches from the opposite end of the hairpin, because the flow
reverses between half-cycles. The melting-temperature cascade is therefore
mirrored with respect to the march index: $\Tm^{\mathrm{dc}}(i) =
\Tm^{\mathrm{ch}}(n_{s}-1-i)$. But $E'$ is measured from a datum of \emph{solid
at the local $\Tm$} (\S\ref{sec:front:enthalpy}), and until v0.5 the state array
was handed to the discharge \textbf{without} being mirrored. The datum therefore
moved by up to \SI{48.9}{\kelvin} across the half-cycle boundary.

The symptom was invisible until $\Tpcm$ was plotted rather than
$\varepsilon_{\mathrm{local}}$: at the start of the discharge the model reported
PCM at \SI{176.6}{\celsius} against a charging inlet of
\SI{160.0}{\celsius} --- hotter than any fluid that had ever touched it. Read
against the charge cascade the same state tops out at exactly
\SI{160.00}{\celsius}, which is correct, and that contrast is what identified
the error.

Two changes fix it. The discharge cascade is now built as the exact reverse of
the charge cascade, which additionally removes a one-layer
(\SI{6.111}{\kelvin}) mismatch that \code{layer\_map} introduces whenever
$n_{s}$ is not a multiple of $N_{\mathrm{lay}}$; and the state array is mirrored
at each half-cycle boundary.

\begin{center}
\begin{tabular}{@{}l r r r@{}}
\toprule
 & before & after & change \\
\midrule
$\Nwells$ (first cycle) & 11.2551 & 11.2551 & --- \\
flow ratio & 1.0262 & 1.0473 & \SI{+2.06}{\percent} \\
$\eta_{\mathrm{RTE}}$ & 0.41475 & 0.41391 & \SI{-0.20}{\percent} \\
$\eta_{\mathrm{storage}}$ & 0.95318 & 0.95266 & \SI{-0.055}{\percent} \\
\midrule
CSS deviation & \SI{-3.34}{\percent} & \SI{-5.74}{\percent} & \\
\bottomrule
\end{tabular}
\end{center}

\noindent
$\Nwells$ is unchanged because the field is sized on the charge, which was never
affected; everything downstream of the discharge moves. The effect on the CSS
result is much larger than on the single cycle, which is what one would expect
of an error that compounds across repeated half-cycle boundaries.
